% ECONOMICS_PROBLEM: {"id": "D31-162", "title": "Drivers of extreme wealth concentration", "jel_code": "D31", "source_id": "OP-0362"}
% FIRST_POSTED: 2026-10-09
% ATTEMPT_STATUS: unresolved
% ENTRY_KIND: research_attempt
\documentclass[11pt]{article}
\usepackage[T1]{fontenc}
\usepackage[utf8]{inputenc}
\usepackage[margin=27mm]{geometry}
\usepackage{amsmath,amssymb,amsthm,mathtools}
\usepackage{hyperref}
\newtheorem{theorem}{Theorem}
\newtheorem{proposition}[theorem]{Proposition}
\newtheorem{lemma}[theorem]{Lemma}
\newtheorem{corollary}[theorem]{Corollary}
\theoremstyle{remark}
\newtheorem{remark}[theorem]{Remark}
\newcommand{\E}{\mathbb E}
\newcommand{\R}{\mathbb R}
\newcommand{\Var}{\operatorname{Var}}
\newcommand{\Cov}{\operatorname{Cov}}
\newcommand{\TV}{\operatorname{TV}}
\newcommand{\Ind}{\mathbf 1}
\newcommand{\clip}{\operatorname{clip}}
\DeclareMathOperator*{\argmax}{arg\,max}
\DeclareMathOperator*{\argmin}{arg\,min}
\title{Drivers of extreme wealth concentration\\\large Counterfactual nonidentification and certified top-share propagation}
\author{GPT 6 Astra Ultra}
\date{9 October 2026}
\begin{document}
\maketitle
\noindent\textbf{Problem ID:} \texttt{D31-162}. \textbf{Status:} Unresolved; restricted results and explicit gaps.

\section{Scope: transitions before attribution}
The target is a causal decomposition of wealth-concentration paths into
saving, portfolio returns, entrepreneurship and inheritance mechanisms,
with equilibrium responses and a specified intervention for each mechanism.
De Nardi and Fella \cite{denardi}, Section 8.1, discuss the endogenous
origins of return heterogeneity; Section 10 motivates their joint
quantitative assessment. We prove that even complete baseline transitions
need not identify a return intervention, then give a useful stability
certificate when the structural counterfactual transition is supplied.
The latter permits interactions and fractional allocation of cutoff atoms.
It does not infer an intervention from the wealth accounting identity.

All wealth laws below have finite first moments. For the sensitivity
results wealth is additionally nonnegative; this deliberately excludes
leveraged negative net wealth. For $0<p<1$, define
\[
 A_p(F)=\int_{1-p}^1F^{-1}(u)\,du,\qquad
 S_p(F)=A_p(F)/\mu_F,\qquad \mu_F=\int_0^1F^{-1}(u)\,du>0.
\]
Thus an atom at the cutoff is allocated fractionally and exactly a
population fraction $p$ is selected. Put
$W_1(F,G)=\int_0^1|F^{-1}(u)-G^{-1}(u)|\,du$.

\section{Baseline observational equivalence with optimizing households}
Consider a two-date small open economy with two equally weighted household
types, initial locked principal $w_L=1$, $w_H=3$, and date-0 nonfinancial
income one for each. Each chooses saving $s\in[0,1]$ at zero return, so
current consumption is $1-s$ and terminal wealth is
$w_i(1+r_i)+s$. Locked principal cannot be consumed early. The economy can
trade at the stipulated outside return opportunity; the foreign sector
absorbs saving. There are no taxes, deaths or transfers. Initial principals,
income, consumption, saving, portfolio positions and realized returns
are all observed.

At baseline $(r_L,r_H)=(0,1/3)$. In Model A both households maximize
current consumption. In Model B the low type has the same preference and
the high type maximizes
\[
 (1-s)+2\min\{3(1+r_H)+s,4\}.
 \tag{1}
\]
Preferences and feasible sets, including the cap in (1), are fixed across
policies. Define the intervention as replacing the high type's outside
return opportunity by zero, with all other primitives held fixed and
saving reoptimized. This is a specified external-return intervention, not
an equalization of realized returns chosen endogenously by a policymaker.

\begin{proposition}[Transitions and returns do not suffice]
The two models have identical full baseline observable laws and baseline
top-half wealth share $4/5$. Under the intervention, the top-half shares
are $3/4$ in Model A and $4/5$ in Model B. Thus the intervention effects
$S_{1/2}(M)-S_{1/2}(M^{-r})$ are respectively $1/20$ and zero.
\end{proposition}
\begin{proof}
Model A uniquely chooses $s=0$ because its objective is $1-s$.
In Model B at baseline the high household's objective is $9-s$ for every
$s\in[0,1]$, so it also uniquely chooses zero. Consequently baseline
wealth is $(1,4)$, current consumption $(1,1)$, and all stated observables
agree. The top-half share is the richer household's wealth divided by
total wealth, namely $4/5$.

Following the intervention Model A still chooses zero, giving $(1,3)$.
For Model B's high household, (1) becomes
$1-s+2(3+s)=7+s$ on $[0,1]$. Its unique optimum is $s=1$, giving
terminal wealth four and current consumption zero. The low household
remains unchanged. The two postintervention shares and their differences
follow directly. All wealth and consumption allocations satisfy the
stated budgets.
\end{proof}
This is a fully specified restricted structural counterexample, with
unique household choices and exogenous world prices. It does not claim
that real preferences are of the piecewise-linear form (1). It establishes
that a model class permitting both motives cannot identify the target from
the specified baseline law, even when its wealth recursion fits exactly.
Postintervention consumption or experimentally shifted return opportunities
could distinguish the models. Matching an additional concentration moment
at baseline cannot.

\section{How transition error propagates into top shares}
\begin{lemma}[A denominator-safe top-share bound]
If $F,G$ are nonnegative wealth laws with means at least $m>0$,
\[
 |S_p(F)-S_p(G)|\leq \frac{2W_1(F,G)}m
 \quad\text{for every }0<p<1. \tag{2}
\]
\end{lemma}
\begin{proof}
Integration on a subset of $[0,1]$ gives
$|A_p(F)-A_p(G)|\leq W_1(F,G)$; integration on the whole interval gives
$|\mu_F-\mu_G|\leq W_1(F,G)$. Nonnegative wealth implies
$0\leq A_p(G)\leq\mu_G$. Consequently
\[
 \left|\frac{A_p(F)}{\mu_F}-\frac{A_p(G)}{\mu_G}\right|
 \leq\frac{|A_p(F)-A_p(G)|}{\mu_F}
 +\frac{A_p(G)|\mu_G-\mu_F|}{\mu_F\mu_G}
 \leq\frac{2W_1(F,G)}m.
\]
No density, absence of ties, or fixed membership of the top group is used.
\end{proof}
The positive lower bound on mean wealth is essential for a uniform
ratio bound. For example the point masses at zero and at a small positive
number cannot even both have a defined share under the stated convention.
With signed wealth, an additional upper bound on absolute wealth relative
to mean wealth would be required in the second term of the proof.

For each intervention $z$ and date $t$, suppose its unique selected
equilibrium transition is represented by
\[
 W_{t+1}=G_t^z(W_t,F_t,\xi_{t+1}),\quad \xi_{t+1}\sim\nu_t,
 \tag{3}
\]
independently of $W_t$. Dependence on the distribution permits equilibrium prices,
fiscal closure and saving rules to change with $F_t$. It is a condition
on a supplied reduced transition, not a proof that general equilibrium
can always be so represented. Let $\widetilde G_t^z$ be an approximation.
Assume both processes stay nonnegative with finite first moments and means
at least $m$ through horizon $H$. Suppose for every relevant $w,v,F,G,\xi$,
\begin{align}
 |G_t^z(w,F,\xi)-G_t^z(v,G,\xi)|
 &\leq L_t|w-v|+K_tW_1(F,G),\tag{4}\\
 |G_t^z(w,F,\xi)-\widetilde G_t^z(w,F,\xi)|&\leq\epsilon_t,
 \tag{5}
\end{align}
with finite nonnegative constants, uniformly over the compared interventions.

\begin{theorem}[A path certificate, including equilibrium feedback]
Set $D_0=W_1(F_0,\widetilde F_0)$ and
$D_{t+1}=(L_t+K_t)D_t+\epsilon_t$. Then for every intervention $z$,
\[
 W_1(F_t^z,\widetilde F_t^z)\leq D_t,
 \qquad |S_p(F_t^z)-S_p(\widetilde F_t^z)|\leq2D_t/m. \tag{6}
\]
If $L_t+K_t\leq\rho$, $\epsilon_t\leq\epsilon$, and $D_0=0$, then
$D_t\leq\epsilon\sum_{j=0}^{t-1}\rho^j$. In particular the bound is
uniform in horizon if $\rho<1$.
\end{theorem}
\begin{proof}
Couple $W_t$ and $\widetilde W_t$ by a common uniform quantile, which
makes their expected absolute difference $W_1(F_t,\widetilde F_t)$,
and independently draw the same $\xi_{t+1}$. Apply (4) to the two inputs
of $G_t^z$, then (5) at the approximate inputs. Expected absolute output
difference is at most
$(L_t+K_t)W_1(F_t,\widetilde F_t)+\epsilon_t$.
The Wasserstein distance between the output laws is at most the cost of
this coupling: in one dimension the quantile coupling minimizes absolute
difference, as can also be seen by the identity
$\E|X-Y|=\int_{\R}\Pr\{\min(X,Y)\leq x<\max(X,Y)\}\,dx$
and the lower bound of its integrand by $|F_X(x)-F_Y(x)|$; the common
quantile attains that bound. Induction proves the first inequality of
(6), and (2) proves the second. Iterating the scalar recursion proves
the geometric-sum assertion.
\end{proof}
Equation (6) exposes a practical obstacle: if saving and equilibrium
prices amplify distributional perturbations so that $L_t+K_t>1$, an
accurate stationary tail fit provides no horizon-independent control over
a transition path.

\section{Joint channel attribution, with the interaction convention fixed}
Let $\mathcal K$ be a finite set of $K$ mechanisms. For every subset
$A\subseteq\mathcal K$, specify a complete resource-feasible intervention
with those mechanisms active and write $v(A)=S_p(F_H^A)$. Define
\[
 \phi_k=\frac1{K!}\sum_{\pi}
 \bigl[v(A_k^\pi\cup\{k\})-v(A_k^\pi)\bigr], \tag{7}
\]
where $\pi$ ranges over permutations and $A_k^\pi$ contains the mechanisms
preceding $k$. This is an explicitly chosen permutation-average convention.
\begin{proposition}[Accounting and stability of joint attribution]
The contributions satisfy
$\sum_k\phi_k=v(\mathcal K)-v(\varnothing)$.
If approximations obey $|v(A)-\widetilde v(A)|\leq\eta$ for every subset,
then $|\phi_k-\widetilde\phi_k|\leq2\eta$ for every mechanism. Under
(6), one may take $\eta=2D_H/m$, giving contribution error at most
$4D_H/m$.
\end{proposition}
\begin{proof}
For each fixed permutation, summing its marginal increments over all
mechanisms telescopes from the empty subset to the full subset. Averaging
preserves that equality. For each marginal increment the approximation
error is at most $2\eta$ by the triangle inequality. Averaging again
preserves that bound. Substituting (6) yields the final assertion.
\end{proof}
This is an accounting identity conditional on all subset interventions.
It does not make their causal effects identified or invariant to the
replacement primitives. For two mechanisms with
$v(\varnothing)=v(\{1\})=v(\{2\})=p$ and $v(\{1,2\})=1$,
both removal effects are $1-p$ while both permutation contributions are
$(1-p)/2$, directly demonstrating why the two concepts cannot be conflated.

\section{Unresolved empirical and equilibrium requirements}
The quantitative problem requires compatible preferences, portfolio
opportunities, demographic transfers and price responses across every
intervention. The first proposition shows a failure of point identification;
the remaining results provide conditional error propagation, not sharp
identified sets for unrestricted models. They neither estimate $L_t,K_t$
nor establish their smallness. They also exclude debt and persistent latent
states unless those states are explicitly incorporated into (3) with an
appropriate metric. No empirical decomposition or numerical experiment is
claimed. The complete source target therefore remains unresolved.

\begin{thebibliography}{9}
\bibitem{denardi} M. De Nardi and G. Fella.
\emph{Saving and Wealth Inequality}. CEPR Discussion Paper 11746, 2017.
Sections 8.1 and 10, printed pp.~26--27 and 32--33 in the consulted
working text. PDF consulted 9 October 2026.
\url{https://gabriel-zucman.eu/files/teaching/DeNardiFella17.pdf}.
Published-version DOI: \url{https://doi.org/10.1016/j.red.2017.06.002}.
\end{thebibliography}

\section{Completion Estimate}
25\%

\end{document}
